\section{Introducción al transporte celular}

La membrana plasmática es una barrera selectiva que regula el paso de sustancias entre el interior y exterior celular.

\section{Transporte pasivo}

El transporte pasivo no requiere energía (ATP). Las moléculas se mueven a favor de su gradiente de concentración.

\subsection{Difusión simple}
\begin{itemize}
    \item Movimiento directo a través de la membrana
    \item Para moléculas pequeñas y liposolubles (O$_2$, CO$_2$, H$_2$O)
    \item No requiere proteínas transportadoras
\end{itemize}

\subsection{Difusión facilitada}
\begin{itemize}
    \item Movimiento a través de proteínas transportadoras
    \item Para moléculas grandes o polares (glucosa, aminoácidos)
    \item Tipos: Canales iónicos y transportadores
\end{itemize}

\subsection{Ósmosis}
\begin{itemize}
    \item Difusión de agua a través de una membrana semipermeable
    \item Depende de la concentración de solutos
    \item Soluciones: Hipotónica, isotónica, hipertónica
\end{itemize}

\section{Transporte activo}

El transporte activo requiere energía (ATP) para mover moléculas en contra de su gradiente de concentración.

\subsection{Transporte activo primario}
\begin{itemize}
    \item Utiliza ATP directamente
    \item Ejemplo: Bomba Na$^+$/K$^+$
    \item Mantiene gradientes iónicos
\end{itemize}

\subsection{Transporte activo secundario}
\begin{itemize}
    \item Utiliza gradientes iónicos como fuente de energía
    \item Cotransporte: Simporte y antiporte
\end{itemize}

\section{Comparación de tipos de transporte}

\begin{table}[h]
\centering
\begin{tabular}{|l|c|c|}
\hline
\textbf{Tipo} & \textbf{Energía} & \textbf{Gradiente} \\
\hline
Difusión simple & No & A favor \\
\hline
Difusión facilitada & No & A favor \\
\hline
Ósmosis & No & A favor (agua) \\
\hline
Transporte activo & Sí & En contra \\
\hline
\end{tabular}
\caption{Comparación de mecanismos de transporte}
\end{table}